% Generated draft: requires mathematical and source review.
\documentclass[11pt]{article}
\usepackage[margin=25mm]{geometry}
\usepackage{fontspec}
\setmainfont{Latin Modern Roman}
\usepackage{amsmath,amssymb,mathtools,unicode-math}
\setmathfont{Latin Modern Math}
\usepackage{enumitem,needspace,xurl,hyperref}
\setlength{\parindent}{0pt}
\setlength{\parskip}{5pt}
\setlength{\emergencystretch}{3em}
\newcommand{\N}{\mathbb{N}}
\newcommand{\Z}{\mathbb{Z}}
\newcommand{\Q}{\mathbb{Q}}
\newcommand{\R}{\mathbb{R}}
\newcommand{\C}{\mathbb{C}}
\newcommand{\F}{\mathbb{F}}
\newcommand{\A}{\mathbb{A}}
\newcommand{\PP}{\mathbb{P}}
\newcommand{\E}{\mathbb{E}}
\newcommand{\eps}{\varepsilon}
\newcommand{\dd}{\,\mathrm{d}}
\begin{document}
\section*{UnsolvedMath ID: 30001226}
\textbf{Source problem code:} catalogue entry
\subsection{Definitions and mathematical statement}
Let \(\mathbb D=\{\zeta\in\C:|\zeta|<1\}\), with coordinate functions \(z,w\) on \(\mathbb D^2\). Define
\[
H^2(\mathbb D^2)=\left\{h(z,w)=\sum_{a,b\geq0}c_{ab}z^aw^b:\ \sum_{a,b\geq0}|c_{ab}|^2<\infty\right\},
\qquad \|h\|^2=\sum_{a,b\geq0}|c_{ab}|^2.
\]
A polynomial is homogeneous if all its nonzero monomials have the same total degree. Interpret a submodule generated by homogeneous polynomials as a closed subspace
\[
M=\overline{\operatorname{span}}\{p q_\lambda:p\in\C[z,w],\ \lambda\in\Lambda\}\subseteq H^2(\mathbb D^2),
\]
where each \(q_\lambda\) is homogeneous; no finiteness assumption on the generating family is supplied. Set \(Q=M^\perp\), the Hilbert-space realization of \(H^2(\mathbb D^2)/M\). For any function \(f\) for which multiplication \(M_f:h\mapsto fh\) is bounded on \(H^2(\mathbb D^2)\), define
\[
S_f=P_QM_f\big|_Q,\qquad [A,B]=AB-BA,\qquad C_f=[S_f^*,S_f].
\]
This defines the compression for bounded multipliers, including polynomials, but the source does not specify the intended class of symbols \(f\).

For a compact operator \(T\) on \(Q\), let \(s_1(T)\geq s_2(T)\geq\cdots\geq0\) be its singular values, padded with zeros when appropriate. The trace class and weak Schatten classes are
\[
\mathcal S_1=\left\{T:\sum_{n=1}^{\infty}s_n(T)<\infty\right\},\qquad
\mathcal S_{p,\infty}=\left\{T:\sup_{n\geq1}n^{1/p}s_n(T)<\infty\right\}\quad(p>0).
\]
To make one common convention explicit, write
\[
\mathcal M_{1,\infty}=\left\{T:\sup_{N\geq1}\frac{\sum_{n=1}^{N}s_n(T)}{\log(1+N)}<\infty\right\}.
\]
This logarithmic ideal is commonly called a Macaev ideal and supports Dixmier traces. The weak Schatten ideal \(\mathcal S_{1,\infty}\) is contained in \(\mathcal M_{1,\infty}\); these definitions must not be identified. The supplied record does not fix its conventions for ``Macaev class'' or ``weak trace class.''

For a positive operator \(T\in\mathcal M_{1,\infty}\), a Dixmier trace associated with an admissible generalized limit \(\omega\) has the form
\[
\operatorname{Tr}_{\omega}(T)=\omega\!\left(\left(\frac{1}{\log(1+N)}\sum_{n=1}^{N}s_n(T)\right)_{N\geq1}\right).
\]
For self-adjoint elements of this ideal, use its linear extension from positive operators. Independence of \(\omega\) is an additional property and is not implied merely by ideal membership.

\textbf{Requested statement.} Classify the homogeneous-generated submodules \(M\), or equivalently the resulting quotient modules \(Q\), for which the specified cross-commutators of the compressed coordinate operators belong to the trace class, the intended Macaev class, or a specified weak Schatten class. Read literally, the supplied clean statement asks for
\[
[A,B]\in\mathcal I\qquad\text{for every }A,B\in\{S_z,S_w,S_z^*,S_w^*\},
\]
for each ideal \(\mathcal I\) under consideration. The frequently used alternative family \([S_i^*,S_j]\), \(i,j\in\{z,w\}\), should be distinguished until the intended meaning of ``cross-commutators'' is established. The weak Schatten exponent \(p\) is not specified in the record.

For the symbols and powers intended by the source, compute the Dixmier traces of powers of \(C_f\) whenever those powers belong to the intended weak trace class. In particular, the positive-integer-power portion of this request is to determine
\[
\operatorname{Tr}_{\omega}(C_f^k),\qquad k\in\{1,2,\ldots\},
\]
whenever \(C_f^k\) lies in the stated trace domain, recording any dependence on \(\omega\). The source does not specify the exponent range, the symbol class, a choice of Dixmier trace, or a measurability hypothesis. Thus the integer-power formulation is an explicitly identified portion of the request, not a resolution of those missing specifications.

The original statement asks only for Dixmier traces of commutator powers and describes quotient modules as generated by homogeneous polynomials. The clean statement assigns homogeneous generators to \(M\) and additionally requests ideal-membership classifications. These differences remain source qualifications requiring mathematical review.
\subsection{Short English statement}
Study compressions of coordinate multiplication to quotients of the bidisk Hardy space by submodules generated by homogeneous polynomials. Classify the relevant commutators by operator-ideal membership and compute Dixmier traces of eligible commutator powers. The symbol class, exponent range, ideal conventions, precise commutator family, and the additional classification target require source clarification.
\subsection{Sources}
\begin{itemize}
\item Supplied local UnsolvedMath registry record, permanent ID \(30001226\), exact source code \texttt{catalogue entry}, title: \emph{Dixmier Traces of Commutators on Bidisk Quotient Modules}. Collection: Oberwolfach Reports Open Problems. The supplied proposed year is 2009; no individual proposer is named.
\item Supplied report citation: \emph{Hilbert Modules and Complex Geometry} (2009), Oberwolfach Reports, pp.\ 1047--1100, DOI \texttt{10.4171/owr/2009/20}. Source URL and literature-source link labelled ``original report'': \url{https://doi.org/10.4171/owr/2009/20}. The supplied background describes these reports as documenting mathematical workshops and their open-problem discussions. No problem-specific page is supplied.
\item Supplied literature-source link labelled ``bidisk trace formulas'': \url{https://research.chalmers.se/en/publication/160925}. The record supplies no author, publication title, theorem number, or precise hypotheses for this reference.
\item Literature assessment dated exactly \texttt{2026-08-21}: the supplied curation labels the problem ``partially solved'' and reports trace-class classifications and trace formulas for important bidisk quotient modules, while describing the general Dixmier trace question as broader. This is supplied dated curation, not independently verified; its status and cited claims require independent verification.
\item The registry labels the statement \texttt{corrected\_verified} and states that it was checked against the official report. This is a supplied verification claim, not an independent verification of the statement or of the relationship between its original and clean formulations.
\end{itemize}
\end{document}
